\section{独立性：从两事件到多个事件}
\label{sec:05-Probability/02-Conditional-Probability/03}

你正在设计一个骰子游戏。玩家下一注赌"点数为偶数"，再下一注赌"点数大于4"。两注独立结算，奖金也独立发放。游戏设计者想知道：这两注之间有关联吗？如果一注中了，另一注的中奖概率会变化吗？

先算基本概率。一枚公平骰子，点数为偶数的概率是 \(P(A)=3/6=1/2\)，点数大于4（即5或6）的概率是 \(P(B)=2/6=1/3\)。两注同时中的情况只有一种——点数为6——所以 \(P(A\cap B)=1/6\)。

这里出现了一个有趣的巧合：
\[
P(A\cap B)=\frac16,\qquad P(A)P(B)=\frac12\cdot\frac13=\frac16.
\]
交集概率恰好等于各自概率的乘积。这意味着什么？

用条件概率的语言重写。\(P(B)=1/3>0\)，可以算
\[
P(A\mid B)=\frac{P(A\cap B)}{P(B)}=\frac{1/6}{1/3}=\frac12=P(A).
\]
知道了点数大于4，偶数的概率还是1/2——和什么都不知道时一样。B的发生没有带来关于A的任何信息。

这就是独立性的直觉：知道一个事件的结果，不改变你对另一个事件的概率判断。

\subsection{不是所有事件对都这么"干净"}

换一对事件试试。还是同一枚骰子，A：点数为偶数，C：点数大于3（即4,5,6）。

\(P(A)=1/2\)，\(P(C)=1/2\)，乘积是 \(1/4\)。但 \(A\cap C=\{4,6\}\)，\(P(A\cap C)=2/6=1/3\)。

\(1/3\neq 1/4\)。乘积公式不成立。用条件概率看：
\[
P(A\mid C)=\frac{P(A\cap C)}{P(C)}=\frac{1/3}{1/2}=\frac23\neq\frac12=P(A).
\]
知道了点数大于3，偶数的概率从1/2跳到了2/3——因为大于3的区域里，偶数（4,6）占了三分之二。C的发生改变了我们对A的判断。

两次尝试，一次成立，一次不成立。区别在哪？关键不在于集合有没有重叠——两对事件都有重叠。区别在于概率分配的方式：在B=\{5,6\}里，偶数和奇数的比例（1:1）恰好和全空间一样；在C=\{4,5,6\}里，偶数和奇数的比例（2:1）偏离了全空间的比例。

用一个夸张的例子把这个观点钉牢。考虑一枚不均匀的骰子，6那一面被灌了铅，概率分配变成：
\[
P(1)=P(2)=\cdots=P(5)=0.1,\quad P(6)=0.5.
\]
还是A="偶数"（即2,4,6），B="大于4"（即5,6）。重新计算：
\[
P(A)=0.1+0.1+0.5=0.7,\quad
P(B)=0.1+0.5=0.6,\quad
P(A\cap B)=P(6)=0.5.
\]
而 \(P(A)P(B)=0.7\times0.6=0.42\neq0.5\)。刚才在公平骰子上独立的两个事件，换了概率权重就不再独立了。

独立性是概率分配的性质，不是集合几何的性质。同一个样本空间，换一组概率权重，原本独立的事件可能变不独立，反之亦然。你不能看一眼集合的文氏图就判断独立性——必须算。

\subsection{从直觉到定义}

把直觉写成定义。当 \(P(B)>0\) 时，"B不改变A的概率"就是
\[
P(A\mid B)=P(A).
\]
代入条件概率公式 \(P(A\mid B)=P(A\cap B)/P(B)\)，两边同乘 \(P(B)\)：
\[
P(A\cap B)=P(A)P(B).
\]
后一式不再需要 \(P(B)>0\) 的条件，即使 \(P(A)=0\) 或 \(P(B)=0\) 也有意义。现代概率论就用这个对称形式定义独立性。

\begin{definition}
事件 \(A\) 与 \(B\) 独立，当且仅当
\[
P(A\cap B)=P(A)P(B).
\]
\end{definition}

这个定义干净、对称、可计算。但它只告诉你乘积是否成立，不附带因果含义，也不承诺"两个事件在物理上互不干扰"。定义是检验工具，不是因果声明。

\subsection{独立不是互斥}

初学者最容易混淆的一对概念就是"独立"和"互斥"。互斥事件满足
\[
P(A\cap B)=0.
\]
若 \(P(A)>0\) 且 \(P(B)>0\)，独立则要求
\[
P(A\cap B)=P(A)P(B)>0.
\]
两个概率为正的事件，不可能既互斥又独立。互斥意味着一个发生会排除另一个——这是强烈的负依赖，恰好和"一个的发生不影响另一个的概率"矛盾。

回到掷骰子。A："点数为偶数"，D："点数为奇数"。它们是互斥的，\(P(A\cap D)=0\)。但 \(P(A)P(D)=1/4>0\)。知道点数为偶数，奇数的概率直接从1/2归零——这大概是你能想象到的最强的"不独立"。

反过来，A与B独立，但它们不是互斥的——6同时属于两个集合。二者的交有正概率，这恰是独立性的特征。

所以：互斥是关于集合是否有交集；独立是关于概率是否按乘积分解。两个维度各自独立。

\subsection{乘积公式好用，但别乱用}

连续抛一枚公平硬币。每次正面的概率是 \(1/2\)。如果各次抛掷独立，前 \(n\) 次全是正面的概率就是
\[
\left(\frac12\right)^n.
\]
当 \(n=5\) 时是 \(1/32\approx3.1\%\)；\(n=10\) 时是 \(1/1024\approx0.1\%\)。数字按指数衰减，直观地说明长串全正面是稀有事件。

但这里的乘积之所以成立，是因为我们在建模时额外假设了独立性。仅凭"实验分成 \(n\) 步"并不自动保证乘积——抽牌不放回也分成多步，却不独立。第一节的不放回抽球正是反例：袋中有3红2白，第一次抽出红球后袋中成分变了，第二次抽到红球的概率从 \(3/5=0.6\) 变成了 \(2/4=0.5\)。联合概率 \(P(\text{两次都红})=(3/5)\times(2/4)=3/10=0.3\)，而 \(P(\text{第一次红})\times P(\text{第二次红})=0.6\times0.6=0.36\neq0.3\)。

独立性是模型赋予的结构，不是多步实验的天然属性。每当你把联合概率写成乘积，背后都有一个或明或暗的独立假设。停下来问一句：这个假设来自对称性论证、来自物理随机化，还是仅仅因为这样算最方便？

\subsection{赌徒谬误：乘积直觉用错了地方}

赌徒谬误是滥用乘积直觉的经典案例。连续出了五次正面，赌徒觉得"反面该来了"——他大概在想：六次全正面的概率只有 \(1/64\approx1.6\%\)，太小了，所以下一次很可能是反面。

推理哪里错了？\(P(\text{六次全正面})=1/64\) 确实是对的。但这是在第一次抛掷之前计算的无条件概率。已经观察到五次正面后，我们不再处于"零信息"状态。此时相关的问题是条件概率：
\[
P(\text{第六次正面}\mid\text{前五次正面}).
\]
若各次抛掷真正独立，硬币没有记忆，这个条件概率仍然是 \(1/2\)。前五次的结果已经是确定的事实，不会"迫使"未来做出补偿。

更具体地，条件概率的计算是：
\[
P(\text{六次全正面}\mid\text{前五次正面})
=\frac{P(\text{六次全正面})}{P(\text{前五次正面})}
=\frac{1/64}{1/32}
=\frac12.
\]
大数定律说的是长期频率——在大量试验之后，正面的比例会接近 \(1/2\)，但这不意味着短期的偏差需要被"纠正"。频率的回归是通过稀释（后来的新试验拉平均值），而不是通过反弹（后来的试验反向补偿）。

\subsection{共享参数制造的伪独立}

另一个对称的错误是把共享参数造成的边缘依赖误解为"变量不独立"。假设你先随机选一枚硬币——它可能有偏置 \(p=0.3\) 或 \(p=0.7\)，等可能——然后连续抛掷。给定硬币的偏置后，各次抛掷独立。

但如果你不知道偏置，观察前两次抛掷的结果。前两次都是正面的概率：
\[
P(\text{HH})=\frac12\times(0.3^2)+\frac12\times(0.7^2)=0.5\times0.09+0.5\times0.49=0.29.
\]
两次正面的边缘概率各为：
\[
P(\text{H})=\frac12\times0.3+\frac12\times0.7=0.5.
\]
若它们是独立的，\(P(\text{HH})\) 应为 \(0.5\times0.5=0.25\)。但实际是 \(0.29>0.25\)——正相关。观察到一次正面，提高了"手上的是高偏置硬币"的信念，从而提高了对下一次正面的预测概率。

这里共享的未知参数 \(p\) 在边缘分布中制造了正相关。这种"给定共同原因后独立、不观测原因时相关"的结构，正是下一节条件独立的主题。

\subsection{两两独立不等于相互独立}

两两独立很好验证：检查每一对即可。但三个事件时，情况更微妙。

考虑三事件 \(A,B,C\)。两两独立只要求
\[
P(A\cap B)=P(A)P(B),\quad
P(A\cap C)=P(A)P(C),\quad
P(B\cap C)=P(B)P(C).
\]
相互独立则要求更高：不仅每一对要独立，对任意非空子集，联合概率都应等于各自概率的乘积。特别是三个事件交集的概率必须也按乘积分解：
\[
P(A\cap B\cap C)=P(A)P(B)P(C).
\]
来看经典反例。独立掷两枚公平硬币。定义三个事件：
\begin{itemize}
\tightlist
\item \(A\)：第一枚为正面；
\item \(B\)：第二枚为正面；
\item \(C\)：两枚结果相同。
\end{itemize}

每个事件的概率都是 \(1/2\)。样本空间四个结果各有概率 \(1/4\)：
\[
\begin{array}{c|c|c}
\omega & \text{概率} & A,B,C\text{ 发生情况} \\
\hline
\text{HH} & 1/4 & A,B,C\text{ 全发生} \\
\text{HT} & 1/4 & A\text{ 发生} \\
\text{TH} & 1/4 & B\text{ 发生} \\
\text{TT} & 1/4 & C\text{ 发生}
\end{array}
\]
检查成对的交：
\begin{itemize}
\tightlist
\item \(A\cap B=\{\text{HH}\}\)，概率 \(1/4 = 1/2\times 1/2\)。通过。
\item \(A\cap C=\{\text{HH}\}\)，概率 \(1/4 = 1/2\times 1/2\)。通过。
\item \(B\cap C=\{\text{HH}\}\)，概率 \(1/4 = 1/2\times 1/2\)。通过。
\end{itemize}
\(A,B,C\) 两两独立。

现在看三个事件的交：\(A\cap B\cap C=\{\text{HH}\}\)，概率 \(1/4\)。但乘积公式预言：
\[
P(A)P(B)P(C)=\frac12\cdot\frac12\cdot\frac12=\frac18.
\]
\(1/4\neq 1/8\)。三者不相互独立。

这个例子的核心教训：任意两对之间看不出关联，不意味着三个事件之间没有高阶约束。C定义在A,B之上（"两枚相同"取决于两枚各自的结果），所以一旦A和B同时发生，C就被决定了。这种高阶依赖完美地躲过了所有成对检查。

\subsection{\(n\) 个事件独立需要多少条件}

把定义推广到 \(n\) 个事件。事件 \(A_1,\ldots,A_n\) 相互独立，是指对任意非空指标集 \(I\subseteq\{1,\ldots,n\}\)，都有
\[
P\biggl(\bigcap_{i\in I}A_i\biggr)=\prod_{i\in I}P(A_i).
\]
仅检查全部 \(n\) 个事件的交不够——中间某个子集可能已经依赖，而全集的乘积恰好巧合成立。仅检查所有成对关系也不够——刚才的反例正是这个陷阱。

完整的条件数量是 \(2^n-n-1\)（所有非空非单点子集）。对于 \(n=10\)，这已是 1013 个等式。没人会逐条验证。

在由明确独立机制构建的模型中，不需要暴力检查。把样本空间定义为乘积空间并令联合概率规定为乘积形式，独立性就被结构保证了。定义告诉我们"什么是独立"，模型结构告诉我们"怎样高效地构造独立"。

\subsection{独立与不相关的区别留到后面}

独立通常推出协方差为零（在方差存在的前提下）：若 \(X,Y\) 独立，则
\[
\E[XY]=\E[X]\E[Y],
\]
从而 \(\Cov(X,Y)=0\)。但反向不成立——协方差为零只排除了线性关联，变量之间仍可能存在精心设计的非线性依赖。

具体反例见\hyperref[sec:05-Probability/04-Expectation-and-Moments/03]{``协方差、相关性与条件波动''}。此刻记住结论就够了：独立是不相关之上的更强要求。独立保证所有阶的依赖都不存在，不相关只保证一阶线性依赖不存在。

下一节打开一个更微妙的局面：两个变量原本相关，在给定第三个变量后独立；或者原本独立，却因为共同结果被观察而变得相关。概率结构不是一成不变的——你观察了什么、条件于什么，会改变变量之间的信息关系。
